\documentclass[a4paper, 12pt]{article}  %set doc type and font size
\usepackage{fontspec}
\usepackage[left=2cm,right=2cm,top=2.2cm,bottom=2.2cm]{geometry}

% math tools
\usepackage{amsmath}
\usepackage{amssymb}
\usepackage{amsfonts}

\setlength{\parindent}{0pt}
% 序列標號
\usepackage{enumerate}

% 導入圖形與表格的封包
\usepackage{graphicx}
\usepackage{booktabs}
\usepackage{tikz}
\usetikzlibrary{arrows.meta, positioning, shapes.geometric}

% set chinese
\usepackage{xeCJK}
\xeCJKsetup{AutoFakeBold=true, AutoFakeSlant=true}
\setmainfont{Times New Roman}
\setCJKmainfont{標楷體}
\XeTeXlinebreaklocale "zh"
\XeTeXlinebreakskip = 0pt plus 1pt

\title{Chapter 3　Assignment \small{Due date: May 11, 2026 (Mon)}}
\author{黃丞暘　611211106}
\date{}

\begin{document}
\maketitle % 顯示標題區塊

\section*{3.1}
\textbf{Sol.}
    
$$
H_0:\mu=\mu_0
\qquad \text{versus} \qquad
H_1:\mu\neq \mu_0
$$

Under \(H_0\),
$$
\bar X_i \sim N\left(\mu_0,\frac{\sigma^2}{m}\right)
$$

Hence,
$$
Z_i=\frac{\bar X_i-\mu_0}{\sigma/\sqrt m}\sim N(0,1)
$$

For a two-sided test with significance level $\alpha$,
$$
P_{H_0}\left(|Z_i|>z_{1-\frac{\alpha}{2}}\right)=\alpha
$$

Therefore, the rejection region is
$$
|Z_i|>z_{1-\frac{\alpha}{2}}
$$
$$
\Rightarrow\left|\frac{\bar X_i-\mu_0}{\sigma/\sqrt m}\right|>z_{1-\frac{\alpha}{2}}
$$
$$
\Rightarrow
\frac{\bar X_i-\mu_0}{\sigma/\sqrt m}>z_{1-\frac{\alpha}{2}}
\quad
\text{or}
\quad
\frac{\bar X_i-\mu_0}{\sigma/\sqrt m}<-z_{1-\frac{\alpha}{2}}
$$

Since \(\sigma/\sqrt m>0\),
$$
\bar X_i-\mu_0>z_{1-\frac{\alpha}{2}}\frac{\sigma}{\sqrt m}
\quad
\text{or}
\quad
\bar X_i-\mu_0<-z_{1-\frac{\alpha}{2}}\frac{\sigma}{\sqrt m}
$$
$$
\Rightarrow
\bar X_i>\mu_0+z_{1-\frac{\alpha}{2}}\frac{\sigma}{\sqrt m}
\quad
\text{or}
\quad
\bar X_i<\mu_0-z_{1-\frac{\alpha}{2}}\frac{\sigma}{\sqrt m}
$$

Let
$$
U=\mu_0+z_{1-\frac{\alpha}{2}}\frac{\sigma}{\sqrt m},
\qquad
L=\mu_0-z_{1-\frac{\alpha}{2}}\frac{\sigma}{\sqrt m}
$$

Therefore,
$$
|Z_i|>z_{1-\frac{\alpha}{2}}
\iff
\text{Reject }H_0
\iff
\bar X_i<L
\quad \text{or} \quad
\bar X_i>U
$$

\section*{3.3}
\textbf{Sol.}
\subsection*{(i)}
Assume
$$
\mu_i=
\begin{cases}
\mu_0, & i<\tau \\
\mu_1, & i\ge \tau
\end{cases}
$$

where
$$
\mu_1=\mu_0+\delta,\qquad \delta\neq 0
$$

From equation \((3.4)\),
$$
\hat{\mu}_0=\bar{\bar X}=\frac{1}{n}\sum_{i=1}^{n}\bar X_i
$$

Since there are \(\tau-1\) subgroups with mean \(\mu_0\), and \(n-\tau+1\) subgroups with mean \(\mu_1\),
$$
E(\bar X_i)=
\begin{cases}
\mu_0, & i<\tau \\
\mu_1, & i\ge \tau
\end{cases}
$$

Therefore,
\begin{align*}
E(\hat{\mu}_0)
&=E\left(\frac{1}{n}\sum_{i=1}^{n}\bar X_i\right) 
=\frac{1}{n}\sum_{i=1}^{n}E(\bar X_i)
=\frac{1}{n}\left[(\tau-1)\mu_0+(n-\tau+1)\mu_1\right] \\
&=\frac{1}{n}\left[(\tau-1)\mu_0+(n-\tau+1)(\mu_0+\delta)\right]
=\frac{1}{n}\left[n\mu_0+(n-\tau+1)\delta\right] \\
&=\mu_0+\frac{n-\tau+1}{n}\delta
\end{align*}

Thus,
$$
\operatorname{Bias}(\hat{\mu}_0)=E(\hat{\mu}_0)-\mu_0=\frac{n-\tau+1}{n}\delta
$$

Since
$$
\delta\neq 0
\quad\text{and}\quad
1<\tau\le n
$$

we have
$$
\frac{n-\tau+1}{n}\delta\neq 0\Rightarrow E(\hat{\mu}_0)\neq \mu_0
$$

Therefore,
$$
\hat{\mu}_0\text{ is not an unbiased estimator of }\mu_0
$$

\subsection*{(ii)}
For the sampling distribution, assume
$$
X_{ij}\sim N(\mu_i,\sigma^2)
$$

Then
$$
\bar X_i\sim N\left(\mu_i,\frac{\sigma^2}{m}\right)
$$

Since $\hat{\mu}_0=\frac{1}{n}\sum_{i=1}^{n}\bar X_i$ and \(\bar X_1,\ldots,\bar X_n\) are independent normal random variables,
$$
\hat{\mu}_0\sim N\left(E(\hat{\mu}_0),Var(\hat{\mu}_0)\right)
$$

Compute
$$
Var(\hat{\mu}_0)=Var\left(\frac{1}{n}\sum_{i=1}^{n}\bar X_i\right)
=\frac{1}{n^2}\sum_{i=1}^{n}Var(\bar X_i)
=\frac{1}{n^2}\sum_{i=1}^{n}\frac{\sigma^2}{m}
=\frac{\sigma^2}{mn}
$$

Therefore,
$$
\hat{\mu}_0\sim N\left(\mu_0+\frac{n-\tau+1}{n}\delta,\frac{\sigma^2}{mn}\right)
$$
    
\section*{3.5}
\textbf{Sol.}
$$
n=10,\quad m=5,\quad \alpha=0.0027\Rightarrow z_{1-\frac{\alpha}{2}}\approx 3
$$

\subsection*{(i)}

\begin{align*}
\bar{\bar X}
&=\frac{1}{10}\sum_{i=1}^{10}\bar X_i
=\frac{35.53+36.01+34.74+35.30+35.80+34.41+35.35+35.19+35.05+35.66}{10} \\
&=35.304
\end{align*}

\begin{align*}
\bar R
&=\frac{1}{10}\sum_{i=1}^{10}R_i
=\frac{2.41+3.24+2.24+3.16+1.06+1.80+1.60+1.42+1.77+1.42}{10} \\
&=2.012
\end{align*}

For subgroup size \(m=5\),
$$
d_1(5)=2.326,\qquad d_2(5)=0.864
$$

The $\bar X$ chart,
$$
U_{\bar X}
=\bar{\bar X}+\frac{z_{1-\frac{\alpha}{2}}}{d_1(m)\sqrt m}\bar R
=35.304+\frac{3}{2.326\sqrt 5}\cdot2.012
\approx 36.465
$$
$$
C_{\bar X}=\bar{\bar X}=35.304
$$
$$
L_{\bar X}
=\bar{\bar X}-\frac{z_{1-\frac{\alpha}{2}}}{d_1(m)\sqrt m}\bar R
=35.304-\frac{3}{2.326\sqrt 5}\cdot2.012
\approx 34.143
$$

Hence,
$$
34.143<\bar X_i<36.465,\qquad i=1,2,\ldots,10
$$

The $R$ chart,
$$
U_R
=\left(1+\frac{z_{1-\frac{\alpha}{2}}d_2(m)}{d_1(m)}\right)\bar R
=\left(1+\frac{3\cdot 0.864}{2.326}\right)\cdot 2.012
\approx 4.254
$$
$$
C_R=\bar R=2.012
$$

$$
L_R
=\left(1-\frac{z_{1-\frac{\alpha}{2}}d_2(m)}{d_1(m)}\right)\bar R
=\left(1-\frac{3\cdot 0.864}{2.326}\right)\cdot 2.012
\approx -0.23\Rightarrow L_R=0
$$

Hence,
$$
0<R_i<4.254,\qquad i=1,2,\ldots,10
$$

In conclusion, the process seems to be in statistical control.

\subsection*{(ii)}

The IC mean is estimated by
$$
\hat{\mu}_0=\bar{\bar X}=35.304
$$

The IC standard deviation is estimated by
$$
\hat{\sigma}
=\frac{\bar R}{d_1(m)}
=\frac{2.012}{2.326}
=0.865
$$

Therefore,
$$
\hat{\mu}_0=35.304,\quad \hat{\sigma}=0.865
$$

\subsection*{(iii)}

When the true process mean is not shifted,
$$
\bar X
\sim
N\left(
\hat{\mu}_0,
\frac{\hat{\sigma}^2}{m}
\right)
$$

The \(\bar X\) chart gives a signal if
$$
\bar X<L_{\bar X}
\quad \text{or} \quad
\bar X>U_{\bar X}
$$

Thus,
\begin{align*}
P(\text{signal})
&=P\left(\bar X<L_{\bar X}\right)+P\left(\bar X>U_{\bar X}\right)\\
&=P\left(Z<\frac{L_{\bar X}-\hat{\mu}_0}{\hat{\sigma}/\sqrt m}\right)+P\left(Z>\frac{U_{\bar X}-\hat{\mu}_0}{\hat{\sigma}/\sqrt m}\right) \\
&=P(Z<-3)+P(Z>3)=2[1-\Phi(3)]=0.0027
\end{align*}

\subsection*{(iv)}

Under the IC process,
$$
ARL_0=\frac{1}{\alpha}=\frac{1}{0.0027}=370.37
$$

For a mean shift of size \(1\),
$$
\mu_1=\mu_0+k\sigma=\mu_0+1\Rightarrow k=\frac{\mu_1-\mu_0}{\hat{\sigma}}=\frac{1}{0.865}=1.156
$$

Then
$$
\bar X\sim N\left(\mu_0+1,\frac{\hat{\sigma}^2}{m}\right)
$$

The probability of no signal after the shift is
\begin{align*}
\beta
&=P\left(\mu_0-z_{1-\frac{\alpha}{2}}\frac{\sigma}{\sqrt m}<\bar X<\mu_0+z_{1-\frac{\alpha}{2}}\frac{\sigma}{\sqrt m}\right) \\
&=P\left(-k\sqrt m-z_{1-\frac{\alpha}{2}}<\frac{\bar X-\mu_1}{\sigma/\sqrt m}<-k\sqrt m+z_{1-\frac{\alpha}{2}}\right) \\
&=\Phi\left(-k\sqrt m+z_{1-\frac{\alpha}{2}}\right)-\Phi\left(-k\sqrt m-z_{1-\frac{\alpha}{2}}\right) \\
&=\Phi\left(-1.156\sqrt 5+3\right)-\Phi\left(-1.156\sqrt 5-3\right) \\
&=\Phi(0.415)-\Phi(-5.585) 
\approx 0.6609
\end{align*}

Therefore,
$$
ARL_1=\frac{1}{1-\beta}=\frac{1}{1-0.6609}=2.949
$$

\section*{3.6}
\textbf{Sol.}

From Exercise 3.5,
$$
\bar{\bar X}=35.304
\quad\text{and}\quad
\bar R=2.012
$$
$$
\bar s
=\frac{1}{10}\sum_{i=1}^{10}s_i
=\frac{0.99+1.28+0.97+1.25+0.43+0.70+0.77+0.53+0.66+0.63}{10}
=0.821
$$

For subgroup size \(m=5\),
$$
\quad d_1(5)=2.326,\quad d_2(5)=0.864,\quad d_3(5)=0.940
$$

\subsection*{(i)}

The \(\bar X\) chart based on \(s\)
$$
U_{\bar X,s}=\bar{\bar X}+\frac{z_{1-\frac{\alpha}{2}}}{d_3(m)\sqrt m}\bar s=35.304+\frac{3}{0.94\sqrt 5}(0.821)\approx 36.476
$$
$$
C_{\bar X}=\bar{\bar X}=35.304
$$
$$
L_{\bar X,s}=\bar{\bar X}-\frac{z_{1-\frac{\alpha}{2}}}{d_3(m)\sqrt m}\bar s=35.304-\frac{3}{0.94\sqrt 5}(0.821)\approx 34.132
$$

Hence,
$$
34.132<\bar X_i<36.476,\quad i=1,\ldots,10
$$

The \(s\) chart,
$$
U_s=\left(1+\frac{z_{1-\frac{\alpha}{2}}\sqrt{1-d_3^2(m)}}{d_3(m)}\right)\bar s=\left(1+\frac{3\sqrt{1-0.94^2}}{0.94}\right)(2.012)\approx 1.715
$$
$$
C_s=\bar s=0.821
$$
$$
L_s=\left(1-\frac{z_{1-\frac{\alpha}{2}}\sqrt{1-d_3^2(m)}}{d_3(m)}\right)\bar s=\left(1-\frac{3\sqrt{1-0.94^2}}{0.94}\right)(0.821)\approx -0.076\Rightarrow L_s=0 \text{ as } s_i\ge0
$$

Hence,
$$
0<s_i<1.715,\quad i=1,\ldots,10
$$

In conclusion, the process seems to be in statistical control.
\\

\subsection*{(ii)}

The \(\bar X\) chart based on $R$,
$$
U_{\bar X,R}=\bar{\bar X}+\frac{z_{1-\frac{\alpha}{2}}}{d_1(m)\sqrt{m}}\bar R=35.304+\frac{3}{2.326\sqrt{5}}(2.012)\Rightarrow 36.465
$$
$$
C_{\bar X,R}=\bar{\bar X}=35.304
$$
$$
L_{\bar X,R}=\bar{\bar X}-\frac{z_{1-\frac{\alpha}{2}}}{d_1(m)\sqrt{m}}\bar R=35.304-\frac{3}{2.326\sqrt{5}}(2.012)\Rightarrow 34.144
$$

Hence,
$$
34.144<\bar X_i<36.465,\quad i=1,\ldots,10
$$

The \(R\) chart,
$$
U_R=\left(1+\frac{z_{1-\frac{\alpha}{2}}d_2(m)}{d_1(m)}\right)\bar R=\left(1+\frac{3\cdot 0.864}{2.326}\right)(2.012)\Rightarrow 4.254
$$
$$
C_R=\bar R=2.012
$$
$$
L_R=\left(1-\frac{z_{1-\frac{\alpha}{2}}d_2(m)}{d_1(m)}\right)\bar R=\left(1-\frac{3\cdot 0.864}{2.326}\right)(2.012)\Rightarrow -0.229\Rightarrow L_R=0 \text{ as } R_i\ge0
$$

Hence,
$$
0<R_i<4.254,\quad i=1,\ldots,10
$$

The two \(\bar X\) charts have limits
$$
\bar X\text{-}s:
\quad
[34.132,\ 36.476]
$$
$$
\bar X\text{-}R:
\quad
[34.144,\ 36.465]
$$

These limits are almost identical.

Both the \(s\) chart and the \(R\) chart show no out-of-control signal.

Both methods suggest that the process is in statistical control

\subsection*{(iii)}

Estimators of \(\sigma\)

The range-based estimator is
$$
\hat\sigma_R=\frac{\bar R}{d_1(m)}=\frac{2.012}{2.326}\approx 0.865
$$

The sample-standard-deviation-based estimator is
$$
\hat\sigma_s=\frac{\bar s}{d_3(m)}=\frac{0.821}{0.94}\approx 0.873
$$

The combined-sample standard deviation is
$$
\hat\sigma_{\mathrm{combined}}=s_{\mathrm{combined}}=\sqrt{\frac{1}{N-1}\sum_{i=1}^{10}\sum_{j=1}^{5}\left(X_{ij}-\bar X_{\mathrm{all}}\right)^2}\approx 0.91
$$

The $R$-based estimator is simple, 
but it uses only the minimum and maximum values in each subgroup.
The \(s\)-based estimator uses all observations inside each subgroup and is usually more efficient.
The combined-sample estimator uses all observations,
it mixes within-subgroup and between-subgroup variation,
it may overestimate.
For SPC,
$$
\hat\sigma_R
\quad\text{and}\quad
\hat\sigma_s
$$
are more appropriate for estimating within-subgroup variation.

\section*{3.7}
\textbf{Sol}

The overall probability of Type I error is
$$
\tilde{\alpha}=1-(1-\alpha)^n
\Rightarrow
1-\tilde{\alpha}=(1-\alpha)^n
\Rightarrow
(1-\tilde{\alpha})^{1/n}=1-\alpha
\Rightarrow
\alpha=1-(1-\tilde{\alpha})^{1/n}
$$

\subsection*{(i)}

When
$$
\tilde{\alpha}=0.001
\text{ and }
n=10
$$
$$
\Rightarrow \alpha=1-(1-0.001)^{1/10}=1-(0.999)^{1/10}\approx 0.0001
$$
$$
\Rightarrow z_{1-\frac{\alpha}{2}}\approx 3.890
$$

When
$$
\tilde{\alpha}=0.001
\text{ and }
n=100
$$
$$\Rightarrow \alpha=1-(1-0.001)^{1/100}=1-(0.999)^{1/100}\approx 0.00001
$$
$$
\Rightarrow z_{1-\frac{\alpha}{2}}\approx 4.417
$$

When $n$ increases from 10 to 100, $\alpha$ must become smaller.
The \(\bar X\) chart based on \(R\) has limits
$
U_{\bar X,R}=\bar{\bar X}+\frac{z_{1-\frac{\alpha}{2}}}{d_1(m)\sqrt m}\bar R
\text{ and }
L_{\bar X,R}=\bar{\bar X}-\frac{z_{1-\frac{\alpha}{2}}}{d_1(m)\sqrt m}\bar R
$
Since $z_{1-\frac{\alpha}{2}}$ is larger when $n=100$, the control limits become wider.

\subsection*{(ii)}

When
$$
\tilde{\alpha}=0.0001
\text{ and }
n=10
$$
$$
\Rightarrow \alpha=1-(1-0.0001)^{1/10}=1-(0.9999)^{1/10}\approx 0.00001
$$
$$
\Rightarrow z_{1-\frac{\alpha}{2}}\approx 4.417
$$

When
$$
\tilde{\alpha}=0.0001
\text{ and }
n=100
$$
$$
\Rightarrow \alpha=1-(1-0.0001)^{1/100}=1-(0.9999)^{1/100}\approx 0.000001
$$
$$
\Rightarrow z_{1-\frac{\alpha}{2}}\approx 4.892
$$

Compared with part \(i\), 
$\tilde{\alpha}$ is reduced from $0.001$ to $0.0001$, 
so the corresponding $\alpha$ also becomes smaller.
The corresponding quantile $z_{1-\frac{\alpha}{2}}$ becomes larger,
so the \(\bar X\) chart limits become wider.
Therefore,
$$
\tilde{\alpha}\downarrow
\quad\Rightarrow\quad
\alpha\downarrow
\quad\Rightarrow\quad
z_{1-\frac{\alpha}{2}}\uparrow
\quad\Rightarrow\quad
\text{wider control limits}
$$

\section*{3.8}
\textbf{Sol}

\subsection*{(i)}

Phase I data:

$\bar X_1,\ldots,\bar X_{10}=80.22,\ 78.31,\ 81.40,\ 78.53,\ 81.32,\ 80.54,\ 77.33,\ 79.24,\ 81.44,\ 77.76$

$s_1,\ldots,s_{10}=3.99,\ 5.35,\ 4.79,\ 4.68,\ 3.79,\ 5.78,\ 3.52,\ 4.51,\ 5.16,\ 4.81$

$$
\bar{\bar X}=\frac{1}{10}\sum_{i=1}^{10}\bar X_i=79.609
\quad\text{and}\quad
\bar s=\frac{1}{10}\sum_{i=1}^{10}s_i=4.638
$$

For subgroup size \(m=5\),
$$
d_3(5)=0.940,\quad z_{1-\frac{\alpha}{2}}=3
$$

The \(\bar X\) chart based on \(s\)
$$
U_{\bar X}=\bar{\bar X}+\frac{z_{1-\frac{\alpha}{2}}}{d_3(m)\sqrt{m}}\bar s=79.609+\frac{3}{0.94\sqrt{5}}(4.638)\approx 86.229
$$
$$
C_{\bar X}=\bar{\bar X}=79.609
$$
$$
L_{\bar X}=\bar{\bar X}-\frac{z_{1-\frac{\alpha}{2}}}{d_3(m)\sqrt{m}}\bar s=79.609-\frac{3}{0.94\sqrt{5}}(4.638)\approx 72.989
$$

Hence,
$$
72.989<\bar X_i<86.229,\quad i=1,\ldots,10
$$

The \(s\) chart
$$
U_s=\left(1+\frac{z_{1-\frac{\alpha}{2}}\sqrt{1-d_3(m)^2}}{d_3(m)}\right)\bar s=\left(1+\frac{3\sqrt{1-0.94^2}}{0.94}\right)(4.638)\approx 9.688$$
$$
C_s=\bar s=4.638
$$
$$
L_s=\max\left\{0,\left(1-\frac{z_{1-\frac{\alpha}{2}}\sqrt{1-d_3(m)^2}}{d_3(m)}\right)\bar s\right\}=0
$$

Hence,
$$
0\le s_i<9.688,\quad i=1,\ldots,10
$$

Therefore, 
the first 10 samples show no out-of-control signal.
The process seems to be in statistical control in Phase I.

Then
$$
\bar X_{11}=79.48,\quad
\bar X_{12}=76.74,\quad
\bar X_{13}=81.12,\quad
\bar X_{14}=86.79
$$
$$
s_{11}=4.56,\quad
s_{12}=7.35,\quad
s_{13}=3.91,\quad
s_{14}=6.31
$$

For the \(\bar X\) chart,
$$
\bar X_{11}, \,X_{12}, X_{13}\,\in[72.989,\ 86.229]
\quad \text{and}\quad
\bar X_{14}\notin[72.989,\ 86.229]
$$

Thus, $\bar X_{14}$ gives an out-of-control signal in the \(\bar X\) chart.

For the \(s\) chart,
$$
0\le s_{11},s_{12},s_{13},s_{14}<9.688
$$

Thus, no signal is found in the \(s\) chart.

\subsection*{(ii)}

For the first 10 samples,
$$
\overline{s^2}=\frac{1}{10}\sum_{i=1}^{10}s_i^2=21.965
$$

Since
$$
\frac{(m-1)s_i^2}{\sigma^2}\sim \chi^2_{m-1}\Rightarrow E(s_i^2)=\sigma^2\Rightarrow Var(s_i^2)=\frac{2\sigma^4}{m-1}\Rightarrow \hat{s_i^2}=\bar {s^2}\sqrt{\frac{2}{m-1}}
$$

So the \(s^2\) chart
$$
U_{s^2}
=\left(1+z_{1-\frac{\alpha}{2}}\sqrt{\frac{2}{m-1}}\right)\overline{s^2}
=\left(1+3\sqrt{\frac{2}{5-1}}\right)(21.965)\approx 68.561
$$
$$
C_{s^2}=\overline{s^2}=21.965
$$
$$
L_{s^2}
=\left(1-z_{1-\frac{\alpha}{2}}\sqrt{\frac{2}{m-1}}\right)\overline{s^2}
=\left(1-3\sqrt{\frac{2}{5-1}}\right)(21.965)\approx -24.631\Rightarrow L_{s^2}=0 \text{ as } s_i^2\ge0
$$

Hence,
$$
0\le s_i^2<68.561,\quad i=1,\ldots,10
$$

The process variability seems to be in statistical control at the first 10 time points

For Phase II,
$
s_{11}^2=4.56^2=20.794
　　\text{and}　　
s_{12}^2=7.35^2=54.023
　　\text{and}　　
s_{13}^2=3.91^2=15.288
　　\text{and}　　
s_{14}^2=6.31^2=39.816
$

All satisfy
$$
0\le s_i^2<68.561,\quad i=11,12,13,14
$$

Therefore,
The $s^2$ chart gives no out-of-control signal in Phase II

\subsection*{(iii)}

Taking the square root of the control limits of the \(s^2\) chart gives a new \(s\) chart
$$
U_s^*
=
\sqrt{U_{s^2}}
=
\sqrt{68.561}
\approx 8.280
$$
$$
C_s^*
=
\sqrt{C_{s^2}}
=
\sqrt{21.965}
\approx 4.687
$$
$$
L_s^*
=
\sqrt{L_{s^2}}
=
0
$$

Thus, the new \(s\) chart is $0\le s_i<8.280$. Square-root $s^2$ chart:[0,\ 8.280]

The original \(s\) chart in part \(i\) is $0\le s_i<9.688$. Original $s$ chart: [0,\ 9.688]

The square-root version has a smaller upper control limit.

For Phase I,
$$
s_i<8.280,\quad i=1,\ldots,10
$$

both \(s\) charts give no out-of-control signal.

For Phase II,
$$
s_{11}=4.56,\quad
s_{12}=7.35,\quad
s_{13}=3.91,\quad
s_{14}=6.31
$$

all satisfy
$$
s_i<8.280,\quad i=11,12,13,14
$$

The square-root version also gives no out-of-control signal in Phase II.

\section*{3.9}
\textbf{Sol}

$$
\bar X=\frac{1}{20}\sum_{i=1}^{20}X_i=24.500
$$

\subsection*{(i)}

For moving window size $\tilde m=2$

$$
d_1(2)=1.128,\quad d_2(2)=0.853
$$

The moving ranges are
$$
MR_i=|X_{i+1}-X_i|,\quad i=1,\ldots,19
$$
$$
MR_i=1,\ 15,\ 13,\ 1,\ 3,\ 2,\ 3,\ 7,\ 4,\ 0,\ 2,\ 6,\ 10,\ 1,\ 1,\ 5,\ 4,\ 2,\ 1
$$
$$
\overline{MR}=\frac{1}{19}\sum_{i=1}^{19}R_i\approx 4.263
$$

The \(\bar X\) chart
$$
U_{\bar X}=\bar X+\frac{z_{1-\frac{\alpha}{2}}}{d_1(m)}\overline{MR}=24.500+\frac{3}{1.128}(4.263)\approx 35.838
$$
$$
C_{\bar X}=\bar X=24.500
$$
$$
L_{\bar X}=\bar X-\frac{z_{1-\frac{\alpha}{2}}}{d_1(m)}\overline{MR}=24.500-\frac{3}{1.128}(4.263)\approx 13.162
$$

Hence
$$
13.162<X_i<35.838,\quad i=1,\ldots,20
$$

But
$$
X_3=39>35.838
$$

Thus the \(\bar X\) chart gives an out-of-control signal at \(X_3\).

The \(R\) chart
$$
U_R=\left(1+z_{1-\frac{\alpha}{2}}\frac{d_2(2)}{d_1(2)}\right)\overline{MR}=\left(1+3\frac{0.853}{1.128}\right)(4.263)\approx 13.935
$$
$$
C_R=\overline{MR}=4.263
$$
$$
L_R=\left(1-z_{1-\frac{\alpha}{2}}\frac{d_2(2)}{d_1(2)}\right)\overline{MR}=\left(1-3\frac{0.853}{1.128}\right)(4.263)\approx -5.408\Rightarrow L_R=0 \text{ as } R_i\ge0
$$

Hence
$$
0\le R_i<13.935,\quad i=1,\ldots,19
$$

But
$$
R_2=|X_3-X_2|=|39-24|=15>13.935
$$

Thus the \(R\) chart gives an out-of-control signal at moving window \(R_2\).

\subsection*{(ii)}

For moving window size $\tilde m=5$
$$
d_1(5)=2.326,\quad d_2(5)=0.864
$$

The moving ranges are
$$
MR_i=\max\{X_i,\ldots,X_{i+4}\}-\min\{X_i,\ldots,X_{i+4}\},\quad i=1,\ldots,16
$$
$$
MR_i=15,\ 17,\ 17,\ 5,\ 7,\ 7,\ 7,\ 7,\ 12,\ 10,\ 10,\ 10,\ 10,\ 5,\ 5,\ 5
$$
$$
\overline{MR}=\frac{1}{16}\sum_{i=1}^{16}R_i\approx 9.313
$$

The \(\bar X\) chart
$$
U_{\bar X}=\bar X+\frac{z_{1-\frac{\alpha}{2}}}{d_1(m)}\overline{MR}=24.500+\frac{3}{2.326}(9.313)\approx 36.511
$$
$$
C_{\bar X}=\bar X=24.500
$$
$$
L_{\bar X}=\bar X-\frac{z_{1-\frac{\alpha}{2}}}{d_1(m)}\overline{MR}=24.500-\frac{3}{2.326}(9.313)\approx 12.489
$$

Hence
$$
12.489<X_i<36.511,\quad i=1,\ldots,20
$$

But
$$
X_3=39>36.511
$$

Thus the \(\bar X\) chart gives an out-of-control signal at \(X_3\).

The \(R\) chart
$$
U_R=\left(1+z_{1-\frac{\alpha}{2}}\frac{d_2(m)}{d_1(m)}\right)\bar R=\left(1+3\frac{0.864}{2.326}\right)(9.313)\approx 19.690
$$
$$
C_R=\bar R=9.313
$$
$$
L_R=\left(1-z_{1-\frac{\alpha}{2}}\frac{d_2(m)}{d_1(m)}\right)\bar R=\left(1-3\frac{0.864}{2.326}\right)(9.313)\approx -1.065\Rightarrow L_R=0 \text{ as } R_i\ge0
$$

Hence
$$
0\le R_i<19.690,\quad i=1,\ldots,16
$$

All moving ranges satisfy
$$
0\le R_i<19.690
$$

Thus the \(R\) chart gives no out-of-control signal.

For \(\tilde m=2\)
$$
\bar X\text{ chart gives a signal at }X_3=39
$$
$$
R\text{ chart gives a signal at }R_2=15
$$

For \(\tilde m=5\)
$$
\bar X\text{ chart gives a signal at }X_3=39
$$
$$
R\text{ chart gives no signal}
$$

The major finding is $X_3=39$,it is an unusually large observation. 
The moving window size \(m=2\) is more sensitive to local jumps between adjacent observations.
The moving window size \(m=5\) produces a larger average moving range and wider \(R\) chart limits.


\section*{3.11}
\textbf{Sol}

\subsection*{(i)}

The sample proportions are
$$
p_i=\frac{X_i}{100},\quad i=1,\ldots,20
$$

The average sample proportion is
$$
\bar p=\frac{\sum_{i=1}^{20}X_i}{20(100)}=\frac{392}{2000}\approx0.196
$$

The \(p\) chart is
$$
U_p=\bar p+z_{1-\frac{\alpha}{2}}\sqrt{\frac{\bar p(1-\bar p)}{m}}=0.196+3\sqrt{\frac{0.196(1-0.196)}{100}}\approx0.315
$$
$$
C_p=\bar p=0.196
$$
$$
L_p=\bar p-z_{1-\frac{\alpha}{2}}\sqrt{\frac{\bar p(1-\bar p)}{m}}=0.196-3\sqrt{\frac{0.196(1-0.196)}{100}}\approx0.077
$$

Hence,
$$
0.077<p_i<0.315
\Rightarrow
7.691<X_i<31.509,\quad i=1,\ldots,20
$$

All samples are within the control limits, the process seems to be in statistical control.

\subsection*{(ii)}

Assume the IC proportion is $\pi_0=0.19$ and the shifted proportion is $\pi_1=0.29$.

Using the \(p\) chart based on formula \((3.19)\),
$$
U_p=\pi_0+z_{1-\frac{\alpha}{2}}\sqrt{\frac{\pi_0(1-\pi_0)}{m}}
$$

To detect the upward shift with probability \(0.9\), require
$$
P_{\pi_1}
\left(
\hat p>U_p
\right)
\ge 0.90
$$

We know that
$$
\hat p\sim N\left(\pi_1,\frac{\pi_1(1-\pi_1)}{m}\right)
$$

Thus,
$$
P_{\pi_1}\left(\hat p>U_p\right)=0.90
\Rightarrow \frac{\pi_1-U_p}{\sqrt{\frac{\pi_1(1-\pi_1)}{m}}}=z_{0.90}
\Rightarrow \frac{\pi_1-\pi_0-z_{1-\frac{\alpha}{2}}\sqrt{\frac{\pi_0(1-\pi_0)}{m}}}{\sqrt{\frac{\pi_1(1-\pi_1)}{m}}}=z_{0.90}
$$
$$
\Rightarrow \sqrt m(\pi_1-\pi_0)=z_{1-\frac{\alpha}{2}}\sqrt{\pi_0(1-\pi_0)}+z_{0.90}\sqrt{\pi_1(1-\pi_1)}
$$
$$
\Rightarrow m=\left[\frac{z_{1-\frac{\alpha}{2}}\sqrt{\pi_0(1-\pi_0)}+z_{0.90}\sqrt{\pi_1(1-\pi_1)}}{\pi_1-\pi_0}\right]^2
$$
$$
\Rightarrow m=\left[\frac{3\sqrt{0.19(1-0.19)}+1.282\sqrt{0.29(1-0.29)}}{0.29-0.19}\right]^2\approx 309.205
$$

Hence, the minimum sample size is $m=310$.

\section*{3.13}
\textbf{Sol}

$$
\sum_{i=1}^{20}X_i=82
\quad\text{and}\quad
\bar p=\frac{\sum_{i=1}^{20}X_i}{nm}=\frac{82}{20(50)}=0.082
$$

\subsection*{(i)}

$$
U_p=\bar p+z_{1-\frac{\alpha}{2}}\sqrt{\frac{\bar p(1-\bar p)}{m}}=0.082+3\sqrt{\frac{0.082(1-0.082)}{50}}\approx 0.198
$$
$$
C_p=\bar p=0.082
$$
$$
L_p=\bar p-z_{1-\frac{\alpha}{2}}\sqrt{\frac{\bar p(1-\bar p)}{m}}=0.082-3\sqrt{\frac{0.082(1-0.082)}{50}}\approx -0.034\Rightarrow L_p=0 \text{ as } p_i\ge 0
$$

Hence
$$
0\le p_i<0.198,\quad i=1,\ldots,20
$$

The process seems to be in statistical control.

\subsection*{(ii)}

Because
$$
m\bar p=50(0.082)=4.100\Rightarrow m\bar p<5
$$

The large-sample condition is not satisfied, so use the binomial distribution directly.

Let
$$
X\sim Binomial(50,0.082)
$$

The exact lower limit is
$$
L_p^*=\frac{\max\{a:P(X\le a)\le \frac{\alpha}{2}\}}{m}=\frac{\max\{a:P(X\le a)\le 0.00135\}}{50}
$$

Since
$$
P(X\le 0)\approx 0.014>0.00135
$$

Then
$$
L_p^*=0
$$

The exact upper limit is

$$
U_p^*=\frac{\min\{a:P(X\ge a)\le \frac{\alpha}{2}\}}{m}=\frac{\min\{a:P(X\ge a)\le 0.00135\}}{50}
$$

Since
$$
P(X\ge 11)\approx 0.002>0.00135
\quad\text{and}\quad
P(X\ge 12)\approx 0.001<0.00135
$$

Then
$$
U_p^*=\frac{12}{50}=0.24
$$

The exact \(p\) chart is
$$
L_p^*=0,\quad C_p=0.082,\quad U_p^*=0.24
$$

Since
$$
0\le p_i<0.24,\quad i=1,\ldots,20
$$

there is no out-of-control signal.

Compared with part \((i)\), 
the exact binomial chart has a wider upper control limit.

\subsection*{(iii)}

The actual Type I error probability is
$$
\alpha'=P(X<mL_p^*)+P(X>mU_p^*)=P(X<0)+P(X>12)=0+P(X\ge 13)\approx 0.000147
$$

\subsection*{(iv)}

Using the control limits in \((3.22)\)

$$
\tilde m=m+z_{1-\frac{\alpha}{2}}^2=50+3^2=59
$$
$$
\tilde p=\frac{m\bar p+\frac{z_{1-\frac{\alpha}{2}}^2}{2}}{\tilde m}=\frac{50(0.082)+\frac{9}{2}}{59}\approx 0.146
$$

The control limit is
$$
U_p=\tilde p+z_{1-\frac{\alpha}{2}}\sqrt{\frac{\tilde p(1-\tilde p)}{\tilde m}}=0.146+3\sqrt{\frac{0.146(1-0.146)}{59}}\approx 0.284
$$
$$
C_p=\tilde p=0.146
$$
$$
L_p=\tilde p-z_{1-\frac{\alpha}{2}}\sqrt{\frac{\tilde p(1-\tilde p)}{\tilde m}}=0.146-3\sqrt{\frac{0.146(1-0.146)}{59}}\approx 0.008
$$

Since
$$
0.008\le p_i<0.284,\quad i=1,\ldots,20
$$

there is no out-of-control signal.

In conclusion, 
All three charts show no out-of-control signal.
The (i) gives the narrowest upper control limit.
The (iv) gives the widest upper control limit.

\section*{3.20}
\textbf{Sol}

The target value is
$$
T=\frac{USL+LSL}{2}
=
\frac{2300+2100}{2}
=
2200
$$

\subsection*{(i)}

The point estimate of \(C_p\) is
$$
\hat C_p
=
\frac{USL-LSL}{6s}
=
\frac{2300-2100}{6(50)}
=
\frac{200}{300}
\approx 0.667
$$

The point estimate of \(C_{pk}\) is
$$
\hat C_{pk}
=
\min\left\{
\frac{USL-\bar X}{3s},
\frac{\bar X-LSL}{3s}
\right\}
=
\min\left\{
\frac{2300-2250}{3(50)},
\frac{2250-2100}{3(50)}
\right\}
=
\min\left\{
\frac{50}{150},
\frac{150}{150}
\right\}
=
0.333
$$

The point estimate of \(C_{pm}\) is
$$
\hat C_{pm}
=
\frac{USL-LSL}
{6\sqrt{s^2+(\bar X-T)^2}}
=
\frac{200}
{6\sqrt{50^2+(2250-2200)^2}}
=
\frac{200}
{6\sqrt{2500+2500}}
\approx 0.471
$$

Between \(C_p\) and \(C_{pk}\), \(C_{pk}\) is more appropriate here.
Because $\bar X=2250\ne T=2200$,
the process mean is not centered at the target,
\(C_p\) only measures spread, while \(C_{pk}\) considers both spread and process centering.

\subsection*{(ii)}

Since
$$
\frac{(n-1)s^2}{\sigma^2}\sim \chi^2_{n-1}
\quad \text{and}\quad
\hat C_p=\frac{USL-LSL}{6s}
$$

We have
$$
P\left(\chi^2_{\alpha/2,n-1}\le\frac{(n-1)s^2}{\sigma^2}\le\chi^2_{1-\alpha/2,n-1}\right)=1-\alpha
\iff
P\left(\hat{C_p}\sqrt{\frac{\chi^2_{\alpha/2,n-1}}{n-1}}\le C_p\le \hat{C_p}\sqrt{\frac{\chi^2_{\alpha/2,n-1}}{n-1}}\right)=1-\alpha
$$
So the \(95\%\) confidence interval for \(C_p\) is
$$
\left[
\hat C_p
\sqrt{
\frac{\chi^2_{0.025,n-1}}{n-1}
},
\quad
\hat C_p
\sqrt{
\frac{\chi^2_{0.975,n-1}}{n-1}
}
\right]
=
\left[
0.667\sqrt{\frac{31.555}{49}},
\quad
0.667\sqrt{\frac{70.222}{49}}
\right]
\approx
[0.535,\ 0.798]
$$

\subsection*{(iii)}

We know $X\sim N(2250,50^2)$ anda defective product occurs when $X<2100\quad\text{or}\quad X>2300$

Thus
\begin{align*}
P(\text{defective})
&=P(X<2100)+P(X>2300)
 =\frac{2100-2250}{50}+\frac{2300-2250}{50} \\
&=P(Z<-3)+P(Z>1)
 =0.001+0.159 \\
&=0.16
\end{align*}

\end{document}